\section*{序論}
10月17日朝までの24時間は、オープンウェイトモデルの性能評価と提供先の拡大、視覚推論を中心とする新たな研究、開発支援エージェントの実用改善が並行して話題になった。前日の発表を受けたMetaのMuse Spark Openへの評価が続く一方、Reflection AIのBeamやElorianAIのGlimpseにも注目が集まった。

研究面では、AnthropicがClaudeを用いてNISTの楕円曲線パラメータの由来を調べたと報告した。製品面ではClaude Codeの応答速度、Codexによる保守作業、Gemini 4 Argonの提供拡大が取り上げられた。政策関連には中国のモデル輸出規制案とする未確認投稿も含まれる。本号はGrokの収集結果を整理した速報であり、公式発表と第三者の評価、窓外の背景情報、未検証の主張を区別して記す。
